\section*{Hoofdstuk \ref{ch:b3:advanced-geochronology} --- Geavanceerde geochronologie}

\begin{solution}{exo:b3:advanced-geochronology:1}
$t_{206} = \ln(1.04500)/\lambda_{238} = \qty{283.8}{\Ma}$; $t_{207} =
\ln(1.3300)/\lambda_{235} = \qty{289.6}{\Ma}$. De ouderdommen verschillen
\qty{6}{\Ma}: de zirkoon is discordant. Zijn $t_{207}$ is de oudste, dus hij ligt
onder de \omterm{def:b3:advanced-geochronology:concordia}{concordia}, aan de kant van het loodverlies.
\end{solution}

\begin{solution}{exo:b3:advanced-geochronology:2}
$J = (e^{5.543\times10^{-10}\times2.8201\times10^7} - 1)/1.3129 =
0.015754/1.3129 = 0.01200$. Dan volgt $t = \ln(1 + 0.01200\times10.000)/\lambda =
\ln 1.12/(5.543\times10^{-10}) = \qty{204.5}{\Ma}$.
\end{solution}

\begin{solution}{exo:b3:advanced-geochronology:3}
$T_{\mathrm{CHUR}} = \ln\bigl(1 + (0.51180 - 0.512638)/(0.1200 -
0.1966)\bigr)/\lambda = \ln(1.010940)/\lambda = \qty{1.67}{\Ga}$;
$T_{\mathrm{DM}} = \ln\bigl(1 + (0.51180 - 0.51315)/(0.1200 - 0.2137)\bigr)/\lambda
= \ln(1.014408)/\lambda = \qty{2.19}{\Ga}$.
\end{solution}

\begin{solution}{exo:b3:advanced-geochronology:4}
Zonder verval geldt $t = N/P_0 = \num{2.0e5}/4.96 = \qty{40.3}{kyr}$. Met verval
$t = -\ln(1 - N\lambda/P_0)/\lambda = -\ln(1 - 0.02014)/\lambda =
\qty{40.7}{kyr}$, \qty{1}{\percent} ouder.
\end{solution}

\begin{solution}{exo:b3:advanced-geochronology:5}
Voor kleine $t$ geldt $x \approx \lambda_{235}t$ en $y \approx \lambda_{238}t$, dus
$y \approx (\lambda_{238}/\lambda_{235})x$, met helling $0.1575$. Als het lood
zeer recent is verloren, is $t_2 \approx 0$ en het \omterm{def:b3:advanced-geochronology:concordia}{onderste snijpunt} is $C(0)$,
de oorsprong. Langs een lijn door de oorsprong is $y/x$ constant, en dus ook
$^{207}\mathrm{Pb}^*/^{206}\mathrm{Pb}^* = x/(137.8\,y)$: elk kristal geeft
de lood--lood-ouderdom van het \omterm{def:b3:advanced-geochronology:concordia}{bovenste snijpunt}, de vormingsouderdom.
\end{solution}

\begin{solution}{exo:b3:advanced-geochronology:6}
Gewichten $1.5625$, $1$, $1.2346$, $2.0408$, $W = 5.838$; $\bar x =
\qty{296.09}{\Ma}$, $\sigma = 1/\sqrt W = \qty{0.41}{\Ma}$. De gewogen
gekwadrateerde afwijkingen tellen op tot $2.32$; met $f = 3$ geldt $\mathrm{MSWD} = 0.77$, onder
$1 + 2\sqrt{2/3} = 2.63$: de ouderdommen komen overeen.
\end{solution}

\begin{solution}{exo:b3:advanced-geochronology:7}
$\lambda\zeta g\rho_d\rho_s/\rho_i = 1.55125\times10^{-10}\times350\times0.5
\times\num{1.10e6}\times0.5 = 0.014931$; $t = \ln(1.014931)/\lambda =
\qty{95.5}{\Ma}$. Als het gesteente snel afkoelde, dateert dit zijn passage door de
gedeeltelijke uitgloeizone van apatiet, ongeveer 60 tot \qty{120}{\celsius}; als het
langzaam afkoelde, is het een gemengde ouderdom, en de sporenlengtes moeten uitsluitsel geven.
\end{solution}

\begin{solution}{exo:b3:advanced-geochronology:8}
Per gram: uranium $20\times10^{-6}/238.05\times N_A = \num{5.06e16}$ atomen, waarvan
\num{5.02e16} uranium-238 en \num{3.64e14} uranium-235; thorium
\num{1.038e17}, $\lambda_{232} = \ln 2/\num{1.40e10} = \qty{4.95e-11}{yr^{-1}}$;
helium $10^{-9}N_A = \num{6.02e14}$. De \omterm{def:b3:advanced-geochronology:cosmogenic}{productiesnelheid} is
$8\lambda_{238}{}^{238}\mathrm U + 7\lambda_{235}{}^{235}\mathrm U +
6\lambda_{232}\mathrm{Th} = (6.23 + 0.25 + 3.08)\times10^7 =
\num{9.57e7}\text{ atomen}\,\unit{g^{-1}.yr^{-1}}$, en $t = \num{6.02e14}/\num{9.57e7} =
\qty{6.29}{Myr}$. Gecorrigeerd: $6.29/0.75 = \qty{8.39}{\Ma}$.
\end{solution}

\begin{solution}{exo:b3:advanced-geochronology:9}
$P_0/\lambda = 4.96/\num{4.994e-7} = \num{9.93e6}\text{ atomen}\,\unit{g^{-1}}$, bereikt tot
\qty{90}{\percent} na $\ln 10/\lambda = \qty{4.6}{Myr}$. Bij
\qty{10}{m.Myr^{-1}} $= \qty{1e-3}{cm.yr^{-1}}$: $\rho\varepsilon/\Lambda =
2.7\times10^{-3}/160 = \qty{1.69e-5}{yr^{-1}}$, $\mu = \qty{1.74e-5}{yr^{-1}}$
en $N = 4.96/\num{1.74e-5} = \num{2.85e5}\text{ atomen}\,\unit{g^{-1}}$, 35 keer minder.
\end{solution}

\begin{solution}{exo:b3:advanced-geochronology:10}
Met $W = w_1 + w_2$ geldt $x_1 - \bar x = w_2(x_1 - x_2)/W$ en $x_2 - \bar x =
-w_1(x_1 - x_2)/W$, dus $\sum w_i(x_i - \bar x)^2 = (x_1 - x_2)^2(w_1w_2^2 +
w_2w_1^2)/W^2 = (x_1 - x_2)^2w_1w_2/W = (x_1 - x_2)^2/(\sigma_1^2 +
\sigma_2^2)$; met $f = 1$ is dit de MSWD. Meten beide dezelfde ouderdom,
$x_1 - x_2$ heeft gemiddelde 0 en variantie $\sigma_1^2 + \sigma_2^2$, dus de
verwachting van het kwadraat gedeeld door die variantie is 1.
\end{solution}

\begin{solution}{exo:b3:advanced-geochronology:11}
$C(2700) = (13.283, 0.52019)$, $C(500) = (0.63628, 0.08065)$. $f = (0.52019 -
0.300)/(0.52019 - 0.08065) = 0.501$, de helft van zijn lood. $x = 0.499\times13.283 +
0.501\times0.63628 = 6.948$. $t_{206} = \qty{1691}{\Ma}$, $t_{207} =
\qty{2105}{\Ma}$; de verhouding $x/(137.8\,y) = 0.1681$ geeft door iteratie een
lood--lood-ouderdom van \qty{2537}{\Ma}. Geen van de drie heeft een betekenis: het zijn
mengsels van de twee gebeurtenissen, en alleen de snijpunten dateren ze. De
lood--lood-ouderdom is een minimum voor de kristallisatie.
\end{solution}

\begin{solution}{exo:b3:advanced-geochronology:12}
Geen erosie: $t = -\ln(1 - \num{1.0e6}\times\num{4.994e-7}/4.96)/\lambda =
\qty{212}{kyr}$. Stationaire toestand: $\rho\varepsilon/\Lambda = P_0/N - \lambda =
\qty{4.46e-6}{yr^{-1}}$, $\varepsilon = \qty{2.64e-4}{cm.yr^{-1}} =
\qty{2.6}{m.Myr^{-1}}$. Voor elke $\varepsilon$ onder deze waarde geldt $N =
P_0(1 - e^{-\mu t})/\mu$ legt één $t$ vast, van \qty{212}{kyr} omhoog tot
oneindig: één meting heeft twee onbekenden. Een tweede nuclide met een andere
halveringstijd, zoals aluminium-26 in hetzelfde kwarts, of de concentratie op
diepte in een boorkern, scheidt ze.
\end{solution}

\begin{solution}{pb:b3:advanced-geochronology:1}
\textbf{1.} $t_{206}$ / $t_{207}$, in \unit{\Ma}: Z1 301.9 / 302.2, Z2 284.1 /
286.5, Z3 265.7 / 268.5, Z4 246.1 / 250.1, Z5 225.1 / 229.4. Een fout van
\qty{0.2}{\percent} op $y$ is ongeveer \qty{0.6}{\Ma}: Z1 is concordant binnen
de fout, de andere zijn discordant, steeds sterker van Z2 tot Z5.

\textbf{2.} De zirkonen verloren een toenemend deel van hun lood bij één latere
gebeurtenis: ze liggen steeds verder naar beneden op de koorde. De koorde ligt onder de
concave \omterm{def:b3:advanced-geochronology:concordia}{concordia}, waar $x$ groter is bij een gegeven $y$: $t_{207} > t_{206}$.

\textbf{3.} Z1: $X = 1/0.04794 = 20.859$, $Y = 0.3467/(137.8\times0.04794) =
0.05248$. Z5: $X = 28.145$, $Y = 0.05178$.

\textbf{4.} Gewichten $1.000$, $1.235$, $1.563$, $1.563$, $2.041$ (in
$10^8$); $\bar x = 0.29402$, $\bar y = 0.040892$, $S_{xx} = \num{7.845e5}$.

\textbf{5.} $b = 0.13260$, $a = \bar y - b\bar x = 0.001905$.

\textbf{6.} Bij \qty{310}{\Ma}, $\mathrm dy/\mathrm dx = 0.1575\,
e^{(\lambda_{238} - \lambda_{235})t} = 0.1218$. De \omterm{def:b3:advanced-geochronology:concordia}{discordia} is slechts
\qty{9}{\percent} steiler dan de \omterm{def:b3:advanced-geochronology:concordia}{concordia}: ze snijdt die onder een kleine hoek,
en het \omterm{def:b3:advanced-geochronology:concordia}{bovenste snijpunt} zal onnauwkeurig zijn.

\textbf{7.} Newton vanuit \qty{310}{\Ma}: $t = \qty{311.2}{\Ma}$.

\textbf{8.} Vanuit \qty{100}{\Ma}: $t = \qty{99.7}{\Ma}$.

\textbf{9.} $(y_i - a - bx_i)/\sigma_i = 0.64$, $-0.91$, $0.35$, $-0.24$,
$0.17$; $\mathrm{MSWD} = 1.435/3 = 0.48$, onder de grens $2.63$ voor $f = 3$:
de zirkonen passen op één rechte.

\textbf{10.} $s_b = 1/\sqrt{\num{7.845e5}} = 0.00113$.

\textbf{11.} Met $b - 2s_b = 0.13034$ is het \omterm{def:b3:advanced-geochronology:concordia}{bovenste snijpunt}
\qty{322.9}{\Ma}; met $b + 2s_b = 0.13485$ is het \qty{302.6}{\Ma}. De ouderdom is
$311 \pm \qty{10}{\Ma}$ bij twee standaardafwijkingen.

\textbf{12.} $f = (x(311.2) - x_i)/(x(311.2) - x(99.7))$: Z1 \qty{4.7}{\percent},
Z5 \qty{41}{\percent}.

\textbf{13.} $J = 0.015754/1.3129 = 0.0119999$, $0.01200$.

\textbf{14.} Stap 1: \qty{150.0}{\Ma}; stap 4: \qty{296.6}{\Ma}; stap 10:
\qty{295.0}{\Ma}.

\textbf{15.} Stappen 4 tot 10: zeven opeenvolgende stappen, \qty{82}{\percent} van het
argon-39, zonder trend. Stappen 1 tot 3 stijgen trapsgewijs.

\textbf{16.} Gelijke fouten: het gemiddelde is $2072.9/7 = \qty{296.1}{\Ma}$, $\sigma =
1.1/\sqrt7 = \qty{0.42}{\Ma}$; de gekwadrateerde afwijkingen tellen op tot $3.83$,
$\mathrm{MSWD} = 3.83/(1.21\times6) = 0.53$, onder $2.15$.

\textbf{17.} $\bar R = \sum f_iR_i = 14.270$, $t = \qty{285.2}{\Ma}$: het
argon dat aan de randen verloren ging en in de eerste stappen vrijkwam, verlaagt het gemiddelde, zoals
het een kalium--argonouderdom zou verlagen.

\textbf{18.} $(700 - 375)/(311.2 - 296.1) = \qty{21.5}{\celsius.Myr^{-1}}$
(met een brede bandbreedte, omdat de zirkoonouderdom bekend is tot op $\pm\qty{10}{\Ma}$).

\textbf{19.} Rond \qty{100}{\Ma} was de graniet begraven onder zijn mesozoïsche
bedekking, en een verhitting of een circulatie van vloeistoffen liet de zirkonen lood verliezen.
De biotiet verloor argon aan zijn randen, wat de trap registreert; zijn
laagste stap, \qty{150}{\Ma}, is een mengsel van oud en teruggezet argon. Het plateau
bleef intact: de biotiet werd niet voldoende verhit om al zijn argon te verliezen.

\textbf{20.} De helling van de hoogte tegen de ouderdom volgens de kleinste
kwadraten is \qty{20.1}{m.Myr^{-1}}.

\textbf{21.} $(65 - 10)/30 = \qty{1.83}{km}$; sinds \qty{43.1}{\Ma}, een gemiddelde van
$1830/43.1 = \qty{42}{m.Myr^{-1}}$.

\textbf{22.} $\rho\varepsilon/\Lambda = P_0/N - \lambda = \num{3.415e-5} -
\num{0.050e-5} = \qty{3.365e-5}{yr^{-1}}$; $\varepsilon = \num{3.365e-5}\times
160/2.7 = \qty{1.99e-3}{cm.yr^{-1}} = \qty{19.9}{m.Myr^{-1}}$.

\textbf{23.} $P_0\Lambda/(\rho N) = \qty{20.2}{m.Myr^{-1}}$, \qty{1.5}{\percent}
meer: bij deze snelheid speelt verval nauwelijks een rol.

\textbf{24.} $\Lambda/(\rho\varepsilon) = 160/(2.7\times\num{1.99e-3}) =
\qty{30}{kyr}$. De graniet steeg ongeveer met \qty{20}{m.Myr^{-1}} tussen 52 en
\qty{43}{\Ma}, sindsdien gemiddeld met \qty{42}{m.Myr^{-1}}, en erodeert
met \qty{20}{m.Myr^{-1}} nu: de ontdaking was op een zeker moment in het
Cenozoïcum sneller.

\textbf{25.} De Varden-graniet kristalliseerde
\textbf{\qty[text-series-to-math=true, separate-uncertainty=true]{311 \pm 10}{\Ma}} geleden (\omterm{def:b3:advanced-geochronology:concordia}{bovenste snijpunt},
bij twee standaardafwijkingen).
\end{solution}
